% ============ 第1章 §1.7 平均算子与差分算子 (原书 pp.23-26) ============

\section{平均算子与差分算子}\label{sec:1.7}
\index{0symbolDk@$D_{k}$}

\index{0symbolEk@$E_{k}$}

\index{C@差分算子}

\index{P@平均算子}










这里同样取 N 为自然数, 考虑长度 $2^{N}$ 的数列 $a = \{a_{j}\}_{j=0}^{2^{N}-1}$.

用 $[a]$ 表示实数 a 的整数部分. 当整数 j 满足 $0 \leq j < 2^{N}$ 时, 对 $k = 0, 1, \ldots, N$, $[2^{-k}j]$ 是整数, 由这个数的定义可知它满足 $2^{k}[2^{-k}j] \leq j < 2^{k}[2^{-k}j] + 2^{k}$. 定义 $E_{k}(a) = \{E_{k}(a)_{j}\}_{j=0}^{2^{N}-1}$ 为
\[
E_{k}(a)_{j} = \frac{1}{2^{k}}\sum_{l = 2^{k}[2^{-k}j]}^{2^{k}[2^{-k}j] + 2^{k} - 1}a_{l}.
\]
也就是说, 将数列 $a = \{a_{j}\}_{j=0}^{2^{N}-1}$ 等分割成长度 $2^{k}$ 的组, 用 $a_{j}$ 所属组的数列的平均替换各 $a_{j}$, 所得的数列记为 $E_{k}(a) = \{E_{k}(a)_{j}\}_{j=0}^{2^{N}-1}$. 另外, 对 $k = 0, 1, \ldots, N-1$, 令 $D_{k}(a) = E_{k}(a) - E_{k+1}(a)$. 对 $k = 0, 1, \ldots, N-1$, 将 $E_{k}, D_{k}, E_{N}$ 视为线性映射. 称 $E_{k}, k = 0, 1, \ldots, N$ 为平均算子, 称 $D_{k}, k = 0, 1, \ldots, N-1$ 为差分算子. 由定义, $E_{0}$ 是恒等算子.

\begin{example}\label{exa:1.14}
对例~\ref{exa:1.13} 中考虑的数列, 求 $E_{1}(a)$, 得
\[
2, 2, 2, 2, 1, 1, 3, 3, 6, 6, 1, 1, 6, 6, 3, 3.
\]
\end{example}

为了能直观理解这里定义的算子的性质, 让我们考虑简单的例子. 对数列 $a = \{a_{j}\}_{j=0}^{2^{N}-1}$, $b = \{b_{j}\}_{j=0}^{2^{N}-1}$, 令 $\langle a, b \rangle = \langle a, b \rangle_{\ell^{2}} = \sum_{j=0}^{2^{N}-1}a_{j}b_{j}$.

\begin{example}\label{exa:1.15}
对由 8 个元组成的数列 $a = \{a_{0}, a_{1}, \ldots, a_{7}\}$ 施行 $E_{1}$, 得
\begin{align*}
E_{1}(a) &= \left\{\frac{a_{0}+a_{1}}{2}, \frac{a_{0}+a_{1}}{2}, \frac{a_{2}+a_{3}}{2}, \frac{a_{2}+a_{3}}{2},\right. \\
&\qquad \left.\frac{a_{4}+a_{5}}{2}, \frac{a_{4}+a_{5}}{2}, \frac{a_{6}+a_{7}}{2}, \frac{a_{6}+a_{7}}{2}\right\}.
\end{align*}
另外, 对原来的数列 a 施行 $E_{1}$ 所得数列 $E_{1}(a)$ 再施行 $E_{2}$, 所得数列 $E_{2} \circ E_{1}(a)$ 为
\begin{align*}
E_{2}(a) &= \left\{\frac{a_{0}+a_{1}+a_{2}+a_{3}}{4}, \frac{a_{0}+a_{1}+a_{2}+a_{3}}{4},\right. \\
&\qquad \left.\frac{a_{0}+a_{1}+a_{2}+a_{3}}{4}, \frac{a_{0}+a_{1}+a_{2}+a_{3}}{4},\right. \\
&\qquad \left.\frac{a_{4}+a_{5}+a_{6}+a_{7}}{4}, \frac{a_{4}+a_{5}+a_{6}+a_{7}}{4},\right. \\
&\qquad \left.\frac{a_{4}+a_{5}+a_{6}+a_{7}}{4}, \frac{a_{4}+a_{5}+a_{6}+a_{7}}{4}\right\}.
\end{align*}
这意味着 $E_{2} \circ E_{1} = E_{2}$. 又, 数列 $b = \{b_{0}, b_{1}, \ldots, b_{7}\}$ 与 $E_{2}(a)$ 的内积 $\langle E_{2}(a), b \rangle$ 为
\begin{align*}
&\frac{(a_{0}+a_{1}+a_{2}+a_{3})(b_{0}+b_{1}+b_{2}+b_{3})}{4} \\
&\qquad + \frac{(a_{4}+a_{5}+a_{6}+a_{7})(b_{4}+b_{5}+b_{6}+b_{7})}{4},
\end{align*}
由于它对 a, b 是对称的, 它也是 a 与 $E_{2}(b)$ 的内积 $\langle a, E_{2}(b) \rangle$. 也就是说, 有 $\langle E_{2}(a), b \rangle = \langle a, E_{2}(b) \rangle$.
\end{example}

\begin{theorem}\label{thm:1.9}
对实数列 $a = \{a_{j}\}_{j=0}^{2^{N}-1}$, $b = \{b_{j}\}_{j=0}^{2^{N}-1}$ 及 $k, l = 0, 1, \ldots, N$, 以下成立.

(1) $E_{k} \circ E_{l} = E_{\max(k,l)}$.

(2) $\langle E_{k}(a), b \rangle = \langle a, E_{k}(b) \rangle$.

(3) $\|a\|_{\ell^{2}}^{2} = \|E_{N}(a)\|_{\ell^{2}}^{2} + \sum_{k=0}^{N-1}\|D_{k}(a)\|_{\ell^{2}}^{2}$.
\end{theorem}

\begin{proof}
关于 (1), 虽然直观上显然, 但为了追求完美, 给出证明.

(1) 给出 $E_{k}$ 的矩阵, 是把 $(i,j)$ 分量为定值 $2^{-k}$ 的 $2^{k}$ 阶方阵 $A_{k}$ 取 $2^{N-k}$ 个直和得到的. 由于 $(A_{k} \oplus A_{k})A_{k+1} = A_{k+1}(A_{k} \oplus A_{k}) = A_{k+1}$, 对 $k = 0, 1, \ldots, N-1$, 有 $E_{k} \circ E_{k+1} = E_{k+1} \circ E_{k} = E_{k+1}$. 因而, 对 $k = 0, 1, \ldots, N-2$, 由于 $E_{k} \circ (E_{k+1} \circ E_{k+2}) = (E_{k} \circ E_{k+1}) \circ E_{k+2}$, $(E_{k+2} \circ E_{k+1}) \circ E_{k} = E_{k+2} \circ (E_{k+1} \circ E_{k})$, 得
\[
E_{k} \circ E_{k+2} = E_{k} \circ (E_{k+1} \circ E_{k+2}) = E_{k+1} \circ E_{k+2} = E_{k+2}
\]
以及
\[
E_{k+2} \circ E_{k} = (E_{k+2} \circ E_{k+1}) \circ E_{k} = E_{k+2} \circ E_{k+1} = E_{k+2}.
\]
反复进行这样的计算, 对一般的 k, l 即得 (1).

(2) 由 (1) 中所考虑的 $A_{k}$ 是对称矩阵可知显然.

(3) 由 $D_{k}\,(k = 0, 1, \ldots, N-1)$ 与 $E_{N}$ 的定义, 分解为
\begin{align*}
\|E_{N}(a)\|_{\ell^{2}}^{2} + \sum_{k=0}^{N-1}\|D_{k}(a)\|_{\ell^{2}}^{2}
&= \|E_{N}(a)\|_{\ell^{2}}^{2} + \sum_{k=0}^{N-1}\|E_{k+1}(a)\|_{\ell^{2}}^{2} \\
&\quad - 2\sum_{k=0}^{N-1}\langle E_{k}(a), E_{k+1}(a)\rangle_{\ell^{2}} + \sum_{k=0}^{N-1}\|E_{k}(a)\|_{\ell^{2}}^{2}.
\end{align*}
由 (1), (2),
\begin{align*}
\langle E_{k}(a), E_{k+1}(a)\rangle_{\ell^{2}} &= \langle E_{k}(a), E_{k+1} \circ E_{k+1}(a)\rangle_{\ell^{2}} \\
&= \langle E_{k+1} \circ E_{k}(a), E_{k+1}(a)\rangle_{\ell^{2}} \\
&= \langle E_{k+1}(a), E_{k+1}(a)\rangle_{\ell^{2}} = \|E_{k+1}(a)\|_{\ell^{2}}^{2},
\end{align*}
所以
\begin{align*}
\|E_{N}(a)\|_{\ell^{2}}^{2} + \sum_{k=0}^{N-1}\|D_{k}(a)\|_{\ell^{2}}^{2}
&= \|E_{N}(a)\|_{\ell^{2}}^{2} + \sum_{k=0}^{N-1}\|E_{k}(a)\|_{\ell^{2}}^{2} - \sum_{k=0}^{N-1}\|E_{k+1}(a)\|_{\ell^{2}}^{2} \\
&= \|E_{0}(a)\|_{\ell^{2}}^{2}.
\end{align*}
由于 $E_{0}(a) = a$, 得到 (3).
\end{proof}

由定理~\ref{thm:1.9} 可知, $\{E_{N}\} \cup \{D_{k}\}_{k=0}^{N-1}$ 是分解 $\ell^{2}$ 范数的算子族. 在本节收尾之际, 将定理~\ref{thm:1.9} 变形为用内积表示的形式. 一般地, 对实向量 $a, b \in \mathbb{R}^{2^{N}}$, 有
\[
\langle a, b \rangle_{\ell^{2}} = \frac{1}{4}\|a + b\|_{\ell^{2}}^{2} - \frac{1}{4}\|a - b\|_{\ell^{2}}^{2}.
\]
由此与定理~\ref{thm:1.9}(3), 得
\begin{equation}\tag{1.15}
\langle a, b \rangle_{\ell^{2}} = \langle E_{N}(a), E_{N}(b)\rangle_{\ell^{2}} + \sum_{k=0}^{N-1}\langle D_{k}(a), D_{k}(b)\rangle_{\ell^{2}}.
\label{eq:1.15}
\end{equation}

\begin{exercise}\label{exer:1.22}
对例~\ref{exa:1.13} 中考虑的数列, 求 $E_{0}(a)$, $E_{2}(a)$, $E_{3}(a)$, $E_{4}(a)$, $D_{0}(a)$, $D_{1}(a)$, $D_{2}(a)$, $D_{3}(a)$.
\end{exercise}

\begin{exercise}\label{exer:1.23}
设 $a = \{a_{j}\}_{j=0}^{7}$ 为实数列. 令
\begin{align*}
A &= \frac{a_{0}+a_{1}+\cdots+a_{7}}{8}, \\
B &= \frac{a_{0}-a_{1}}{2}, \quad C = \frac{a_{2}-a_{3}}{2}, \quad D = \frac{a_{4}-a_{5}}{2}, \quad E = \frac{a_{6}-a_{7}}{2}, \\
F &= \frac{a_{0}+a_{1}-a_{2}-a_{3}}{4}, \quad G = \frac{a_{4}+a_{5}-a_{6}-a_{7}}{4}, \\
H &= \frac{a_{0}+a_{1}+a_{2}+a_{3}-a_{4}-a_{5}-a_{6}-a_{7}}{8}.
\end{align*}
计算 $8A^{2} + 2B^{2} + 2C^{2} + 2D^{2} + 2E^{2} + 4F^{2} + 4G^{2} + 8H^{2}$.
\end{exercise}
